\SourceLeaf{271}%
\noindent folgt. Die Constante ist jedoch gleich Null zu setzen, da $y$ für $x = 0$ verschwindet, sodass schliesslich

\[
y = \alpha \,.\, \frac{\sin kx}{k}
\]

wird. Demnach ist:

\[
f(x) = \alpha\, \frac{\sin kx}{k}
\]

und nach 14):

\[
\tag*{(G)} g(x) = \cos kx.
\]

17. Setzt man für $g(x)$ den Ausdruck $\cos kx$ in den Relationen V. 1), 2), 3) und 4) ein, so ergeben sich daraus die 10 Relationen:

\[
\cos kc = \cos ka \,.\, \cos kb,
\]

\[
\begin{array}{cc}
\mathrm{tg}\, kb = \mathrm{tg}\, kc \,.\, \cos A, & \mathrm{tg}\, ka = \mathrm{tg}\, kc \,.\, \cos B, \\[4pt]
\mathrm{tg}\, ka = \mathrm{tg}\, A \,.\, \sin kb, & \mathrm{tg}\, kb = \mathrm{tg}\, B \,.\, \sin ka, \\[4pt]
\sin ka = \sin A \,.\, \sin kc, & \sin kb = \sin B \,.\, \sin kc, \\[4pt]
\cos A = \cos ka \,.\, \sin B, & \cos B = \cos kb \,.\, \sin A,
\end{array}
\]

\[
\cos kc = \cot A \,.\, \cot B,
\]

und das sind die bekannten Gleichungen der \textls{sphärischen Trigonometrie}, wenn der Radius der Kugel gleich $\dfrac{1}{k}$ gesetzt wird, die sich in die Gleichungen der \textls{Nicht-Euklidischen Trigonometrie} verwandeln, wenn man für $k$ einen rein imaginären Werth $k = ik'$ wählt.

Es ist auffallend, dass \textsc{Gauss} die Constante, die bei der Integration der Differentialgleichung 15) auftritt, mit $k$ bezeichnet; ob er dadurch einen Zusammenhang mit dem Krümmungsmass andeuten wollte, das er in den Disquisitiones generales circa superficies curvas ebenfalls $k$ nennt, lässt sich freilich nicht mit Bestimmtheit behaupten. Bemerkt zu werden verdient bei dieser Gelegenheit eine Äusserung von \textsc{Minding} in seiner Abhandlung: \emph{Beiträge zur Theorie der kürzesten Linien auf krummen Flächen}, die 1840 in \textsc{Crelles} Journal für die reine und angewandte Mathematik (Bd. 20. S. 323--327) erschienen war:

»Dass auf jeder Fläche von unveränderlichem positivem Krümmungsmasse zwischen den Seiten und Winkeln eines aus kürzesten Linien gebildeten Dreiecks die Formeln der sphärischen Trigonometrie gelten, folgt sogleich, wenn man sich erinnert, dass jede Fläche dieser Art sich auf eine Kugel abwickeln lässt. Ist das Krümmungsmass negativ, so gelten dieselben Formeln mit der Änderung, dass die hyperbolischen Functionen der Seiten an die Stelle der trigonometrischen treten. Sind nemlich $a, b, c$ die Seiten des Dreiecks, $A$ der Gegenwinkel von $a$ und $k$ das unveränderliche Krümmungsmass, gleichviel ob positiv oder negativ, so ist es nicht schwer, die Richtigkeit folgender Gleichung zu beweisen:

\[
\cos a\BareRoot{k} = \cos b\BareRoot{k} \,.\, \cos c\BareRoot{k} + \sin b\BareRoot{k} \,.\, \sin c\BareRoot{k} \,.\, \cos A.\text{«}
\]

Die Rotationsflächen von constantem negativen Krümmungsmasse hatte \textsc{Minding} bereits 1839 bestimmt (\textsc{Crelles} Journal, Bd. 19). Man vergleiche dazu folgende Notiz von \textsc{Gauss}, die sich in einem Handbuche findet und aus der Zeit zwischen 1823 und 1827 stammt:
